\documentclass[10pt,a4paper]{amsart}
\usepackage[margin=19mm]{geometry}
\usepackage{amsmath,amssymb,mathtools}
\usepackage[scr=rsfs,cal=euler]{mathalfa}
\usepackage{xfrac}
\usepackage[unicode,hidelinks,bookmarks=false]{hyperref}
\newcommand{\ie}{i.e.\ }
\newcommand{\cf}{cf.\ }
\newcommand{\resp}{respectively\ }
\newcommand{\Subsection}{\subsection}
\providecommand{\nobreakdash}{\nobreak\hbox{-}}
\setlength{\emergencystretch}{2em}
\multlinegap0pt
\usepackage{bm}
\usepackage{bbm}
\usepackage{dsfont}
\usepackage{xspace}
\usepackage{cancel}
\usepackage{mathtools}
\usepackage{amsthm}
\makeatletter
\@ifundefined{thm}{\newtheorem{thm}{Theorem}}{}
\@ifundefined{lem}{\newtheorem{lem}[thm]{Lemma}}{}
\makeatother
\newcommand{\AG}[2]{\mathrm{AG}(#1,\,#2)}
\newcommand{\GF}[1]{\mathrm{GF}(#1)}
\newcommand{\PGv}[1]{\mathrm{PG}(#1)}
\newcommand{\PG}[2]{\mathrm{PG}(#1,\,#2)}
\newcommand{\SSS}{\mathcal{S}}
\newcommand{\la}{\lambda}
\title[On the Maximality of Additive Codes]{On the Maximality of Additive Codes}
\author{Tim Alderson}
\date{Source: arXiv:2607.22297v1 (CC BY 4.0)}
\begin{document}
\maketitle

\noindent\emph{Source and licence.} On the Maximality of Additive Codes. Version arXiv:2607.22297v1 (CC BY 4.0), distributed under the Creative Commons Attribution 4.0 International licence, as recorded in the arXiv licence metadata for this exact version and shown on the abstract page https://arxiv.org/abs/2607.22297v1. Full rights notice: licenses/arxiv-2607.22297-CC-BY-4.0.txt.\par\smallskip

\noindent\textit{Editorial scope.} Counted unit: the Thm whose source label is \texttt{thm:dir16} in the article (source0.tex lines 396--398, source label \texttt{thm:dir16}), delivered together with the article's own complete original proof of it (source0.tex lines 445--513, one \texttt{proof} environment of 69 lines). No mathematics is written, repaired or re-derived: statement and proof are the article's own text line for line. The proof is detached from the statement in the source: the article states the result at lines 396--398 and prints its proof at lines 445--513, and the proof environment's own header names the statement, so the association is the article's own. Required context retained uncounted: lem:BC (lines 431--437). Every definition, display and label of those lines is retained unchanged. Omitted from this package and not redistributed: 1--395, 399--430, 438--444, 514--853. Nothing inside a retained excerpt is omitted or altered: every retained body is delivered exactly as the article's lines are written, so no omission receipt of any kind accompanies this package. Citations: the retained excerpts carry no citation command at all, so no bibliography entry of the article is transcribed and no reference is redistributed. Reference adaptation: none was needed and none was made; every reference inside the delivered unit resolves inside it, so no unresolved reference or citation is produced. Printed slips kept literal: no printed word, symbol, hypothesis or number of the counted statement or of its proof was altered. Layout and class: the article's own class is \texttt{article} with packages including amsmath, amssymb, amsthm; the wrapper is the lane's pinned \texttt{amsart} editorial wrapper at 10pt on A4, which restates the theorem-like environments the retained text needs and adds the article's own macro definitions that the retained text needs (\texttt{\textbackslash AG}, \texttt{\textbackslash GF}, \texttt{\textbackslash PG}, \texttt{\textbackslash PGv}, \texttt{\textbackslash SSS}, \texttt{\textbackslash la}), with extra wrapper packages (bm, bbm, dsfont, xspace, cancel, mathtools). The delivered numbering is the unit's own. Assets: no retained span contains a figure environment, a graphics-inclusion command, a PDF-page inclusion, a table or a TikZ picture; no image file is copied into this package, the package declares no asset dependency and no image file is redistributed with it. Licence: version arXiv:2607.22297v1 of this article is distributed under the Creative Commons Attribution 4.0 International licence, as recorded in the arXiv licence metadata for this exact version and shown on the abstract page https://arxiv.org/abs/2607.22297v1. A full rights notice is delivered at licenses/arxiv-2607.22297-CC-BY-4.0.txt. Characters: every retained excerpt is pure ASCII, so the delivered engine, XeLaTeX with its default Latin Modern text font, reports no missing character.
\begin{lem}\label{lem:BC}
Let $B\subseteq E\subseteq\AG52$ with $|B|+|E|=16$ and $6\le|B|\le8$, and let
\[
D(B,E)\;=\;\bigl\{[\beta+\gamma] : \beta\in B,\ \gamma\in E,\ \beta\neq\gamma\bigr\}.
\]
Then $D(B,E)$ contains a plane of $\PG42$.
\end{lem}

\begin{thm}\label{thm:dir16}
Let $T$ be a set of $16$ points of $\AG62$.  Then $D(T)$ contains a solid of $\Pi=\PG52$.
\end{thm}

\begin{proof}[Proof of Theorem~\ref{thm:dir16}]
%\textit{Step 1: $T$ holds a plane.}
The $\binom{16}2=120$ pairs of points of $T$ determine directions among the $63$ points of $\Pi$, so two distinct pairs share a direction and their four points form a $2$-flat $\alpha\subseteq T$.  Let $W\le\GF2^6$ be the $2$-dimensional subspace with $\alpha=v+W$ (equivalently, $W=\{a+b:a,b\in\alpha\}$), i.e.\ the translation subspace of $\alpha$, and let $\ell=\PGv{W}$ be its line at infinity.  Every nonzero $w\in W$ is a difference of two points of $\alpha\subseteq T$, so
\begin{equation}\label{eq:ellD}
\ell\subseteq D(T).
\end{equation}

\smallskip
%\textit{Step 2: the quotient modulo $W$.}
Write $\overline V=\GF2^6/W\cong\GF2^4$ and $\la\mapsto\bar\la$ for the quotient map.  The cosets of $W$ are the $2$-flats parallel to $\alpha$. For $X\in\overline V$ let $\alpha_X$ denote the corresponding $2$-flat and $n_X=|T\cap\alpha_X|$, and let $O=\bar\alpha$, so that
\begin{equation}\label{eq:twelve}
n_O=4,\qquad \sum_{X\neq O}n_X=12 .
\end{equation}
We regard $\overline V$ as the affine space $\AG42$ with hyperplane at infinity $\overline\Pi=\PGv{\overline V}\cong\PG32$.  For a point $x=[\bar v]\in\overline\Pi$, the plane $\gamma_x=\PGv{W+\langle v\rangle}$ of $\Pi$ contains $\ell$, and the four points of $\gamma_x\setminus\ell$ are the $[v+w]$, $w\in W$.  Geometrically, $\overline V$ together with $\overline\Pi$ is the quotient of $\Sigma$ at $\ell$: the points of $\overline\Pi$ are the planes of $\Pi$ through $\ell$ (the plane corresponding to $x$ being $\gamma_x$), and the affine points are the planes of $\Sigma$ meeting $\Pi$ precisely in $\ell$, that is, the $2$-flats parallel to $\alpha$.  We shall call $x$ \textit{covered} if $\gamma_x\setminus\ell\subseteq D(T)$.

\smallskip
\noindent\textit{Claim 1:}  \textit{If all three points of some line of $\overline\Pi$ are covered, then $D(T)$ contains a solid.}\\ 
Indeed, such a line is $\PGv{\overline U}$ for a $2$-dimensional $\overline U\le\overline V$.  Its preimage $U\le\GF2^6$ is $4$-dimensional and contains $W$, and $\PGv U$ is the union of $\ell$ and the three sets $\gamma_x\setminus\ell$, $x\in\PGv{\overline U}$ ($3+3\cdot4=15$ points).  By~\eqref{eq:ellD} and the covered property, the solid $\PGv U$ lies in $D(T)$.

\smallskip
By Claim~1 we may assume for the remainder of the proof that \emph{no line of $\overline\Pi$ has all three of its points covered}.

\smallskip
\noindent\textit{Claim 2.  If $X\neq O$ and $n_X>0$, then $[O+X]$ is covered.}\\
To see this, fix $\la\in T\cap\alpha_X$.  As $\mu$ varies over the four points of $\alpha\subseteq T$, the differences $\la+\mu$ range over a coset of $W$ disjoint from $W$, namely over the four vectors representing the points of $\gamma_x\setminus\ell$ with $x=[O+X]$.  Hence $\gamma_x\setminus\ell\subseteq D(T)$.

\smallskip
Let $\SSS=\{[O+X] : X\neq O,\ n_X>0\}\subseteq\overline\Pi$ and $s=|\SSS|$.  Since $X\mapsto[O+X]$ is a bijection from $\overline V\setminus\{O\}$ to $\overline\Pi$, we have $s=\#\{X\neq O : n_X>0\}$, and by Claim~2 every point of $\SSS$ is
covered.

\smallskip
\noindent\textit{Claim 3.  $n_X+n_Y\le4$ for all distinct $X,Y\neq O$ with $n_X,n_Y\ge1$.}

Indeed, suppose $n_X+n_Y\ge5$.  The points $x=[O+X]$, $y=[O+Y]$ and $z=[X+Y]$ are distinct and collinear in $\overline\Pi$ (their representing vectors sum to $0$).  Let $[u]$ be any of the four points of $\gamma_z\setminus\ell$, so $\bar u=X+Y$.  Translation by $u$ maps $\alpha_X$ bijectively onto $\alpha_Y$. Since $n_X+n_Y>4=|\alpha_Y|$, the sets $(T\cap\alpha_X)+u$ and $T\cap\alpha_Y$ meet, giving $\la\in T\cap\alpha_X$ with $\la+u\in T\cap\alpha_Y$, whence $[u]\in D(T)$.  Thus $z$ is covered, while $x$ and $y$ are covered by Claim~2, so the line $\{x,y,z\}$ is fully covered, contrary to the standing assumption.

\smallskip
We record three consequences.
\begin{itemize}
\item[(a)] \textit{$\SSS$ is a cap of $\overline\Pi$} (no three points collinear): a line of $\overline\Pi$ with all three points in $\SSS$ would consist of covered points.  Caps of $\PG32$ have at most $8$ points (through any point of a cap pass $7$ lines, each containing at most one further cap point), so $s\le8$.
\item[(b)] \textit{$n_X\in\{1,2\}$ whenever $X\neq O$ and $n_X>0$.} Note first that $s\ge3$, by~\eqref{eq:twelve} and $n_X\le4$.  If some $n_X=4$, then any second class $Y$ with $n_Y\ge1$ violates Claim~3.  If some $n_X=3$, then Claim~3 forces $n_Y=1$ for every other class counted by $s$, so $12=3+(s-1)\le10$ by (a), a contradiction. \item[(c)] Let $\kappa=\#\{X : n_X=2\}$.  By~\eqref{eq:twelve} and (b), $2\kappa+(s-\kappa)=12$, so $\kappa=12-s\ge0$ and $s-\kappa=2s-12\ge0$; with (a), 
\begin{equation}\label{eq:srange}
6\le s\le 8,\qquad \kappa=12-s\in\{4,5,6\}.
\end{equation}
\end{itemize}

\smallskip
%\textit{Step 3: common direction on $\ell$, and the quotient.}
For each of the $\kappa$ classes $X$ with $n_X=2$, the set $L_X=T\cap\alpha_X$ is an affine line whose point at infinity lies on $\ell$.  Suppose two of them, $L_X=\{\la_1,\la_2\}$ and $L_Y=\{\mu_1,\mu_2\}$, had distinct points at infinity, i.e.\ $\la_1+\la_2\neq\mu_1+\mu_2$.  The four cross differences $\la_i+\mu_j$ are then pairwise distinct ($\la_1+\mu_1=\la_1+\mu_2$ would give $\mu_1=\mu_2$, and $\la_1+\mu_1=\la_2+\mu_2$ would give
$\la_1+\la_2=\mu_1+\mu_2$), and they represent points of $\gamma_z\setminus\ell$ for $z=[X+Y]$. Since $|\gamma_z\setminus\ell|=4$, the point $z$ is covered, and with Claim~2 the line $\{[O+X],[O+Y],z\}$ is fully covered, contrary to the standing assumption.  Hence all $\kappa$ lines $L_X$ have a common point at infinity $P=[p]\in\ell$.

Write $\widehat V=\GF2^6/\langle p\rangle\cong\GF2^5$ and $\la\mapsto\widehat\la$ for the quotient map, and regard $\widehat V$ as $\AG52$ with hyperplane at infinity $\widehat\Pi=\PGv{\widehat V}\cong\PG42$.  Call an affine line with point at infinity $P$ a \textit{$P$-line}; the fibres of the quotient over the points of $\AG52$ are exactly the $P$-lines, each $P$-line lies in a single class $\alpha_X$ (as $p\in W$), and each class
contains exactly two $P$-lines.  Geometrically, $\widehat V$ together with $\widehat\Pi$ is the quotient of $\Sigma$ at $P$, the affine points being the $P$-lines.  Put
\[
B=\{\widehat\la : \{\la,\la+p\}\subseteq T\},\qquad E=\{\widehat\la : \{\la,\la+p\}\cap T\neq\emptyset\},
\]
so $B\subseteq E\subseteq\AG52$, distinct $P$-lines giving distinct points of $\AG52$.

The two $P$-lines of $\alpha$ lie in $T$ and contribute two points to $B$, and each class $X$ with $n_X=2$ contributes the point of its line $L_X$. No other $P$-line lies in $T$, since a class with $n_X\le1$ contains no two points of $T$, while in a class with $n_X=2$ the second $P$-line misses $T$.  Hence $|B|=2+\kappa=14-s$.  Likewise, a class $X\neq O$ with $n_X>0$ meets $T$ in a single point ($n_X=1$) or
in the single $P$-line $L_X$ ($n_X=2$), so it contributes exactly one point to $E$, and $\alpha$ contributes two: $|E|=2+s$.  By~\eqref{eq:srange},
\[
|B|+|E|=16,\qquad 6\le|B|\le8,\qquad B\subseteq E .
\]

\smallskip
%\textit{Step 4: lifting a plane to a solid.}
By Lemma~\ref{lem:BC} there is a plane $\tau\subseteq D(B,E)\subseteq\widehat\Pi$. Write $\tau=\PGv{\widehat U}$ with $\widehat U\le\widehat V$ of dimension $3$, and let $U\le\GF2^6$ be the preimage of $\widehat U$, of dimension $4$ and containing $p$.  We claim that the solid $\PGv U$ lies in $D(T)$.

First, $P=[p]\in\ell\subseteq D(T)$ by~\eqref{eq:ellD}.  Every other point of $\PGv U$ is $[u]$ with $\widehat u\neq0$ and $R=[\widehat u]\in\tau$, and the points of $\PGv U$ lying over $R$ are exactly the pair $[u]$, $[u+p]$.  Since $\tau\subseteq D(B,E)$, there are $\beta\in B$ and $\gamma\in E$ with $\beta\neq\gamma$ and $R=[\beta+\gamma]$.  Choose $\la$ with $\widehat\la=\beta$, so that $\{\la,\la+p\}\subseteq T$, and a point $\mu\in T$ with $\widehat\mu=\gamma$.  The vectors $\la+\mu$ and $\la+\mu+p$ are nonzero (as $\beta\neq\gamma$), differ by $p$, and reduce to $\beta+\gamma$ modulo $\langle p\rangle$. Therefore $\{[\la+\mu],[\la+\mu+p]\}$ is exactly the pair of points of $\PGv U$ over $R$, and both are differences of points of $T$.  As $R$ ranges over the seven points of $\tau$, this accounts for all $14$ points of $\PGv U\setminus\{P\}$, proving the claim and the theorem. 
\end{proof}

\end{document}
